Angesichts des bedeutenden Plus, mit dem die Bilanzrechnung abschließt, erscheint eine Erhöhung der Pension für die bevorstehende Periode zulässig. Über die Höhe der Erhöhung habe ich Folgendes zu bemerken.

Fragt man, um wie viel der variable Teil der Witwenpension unter Beibehaltung aller übrigen Einrichtungen erhöht werden muss, damit Soll und Haben in der Bilanz vollkommen gleich werden, so ergibt eine einfache Rechnung diese Erhöhung \(=\) 13 Thl. 7 Ggr. Wählt man eine kleinere Erhöhung, so schließt die veränderte Bilanz noch immer mit einem Plus, bei einer größeren Erhöhung dagegen mit einem Minus. Würde also der variable Teil der Pension von jetzt an auf 60 Thaler festgesetzt, sodass jede berechtigte Witwe zusammen 260 Thaler erhielte, solange ihre Anzahl 18 nicht überschreitet, andernfalls aber neben dem festen Teil von 200 Thaler noch den entsprechenden Anteil an der Gesamtsumme 1080 Thaler, so würde die auf dieselbe Art wie oben geführte Bilanzrechnung noch mit einem Plus von 1695 Thalern abschließen. Man ist zu dieser Maßnahme also nicht nur vollkommen berechtigt, sondern hat auch die Aussicht, dass nach wenigen Jahren eine weitere Erhöhung möglich sein wird, sofern keine \textit{großen} Verluste eintreten, der Vorteil des höheren Zinsfußes fortbesteht und bei der jetzt noch nicht erreichten Witwenzahl 18 von der in Rechnung gestellten jährlichen Summe vorerst jedes Jahr etwas übrig bleibt.

Mit einem Minus von 877 Thalern würde die Bilanz dagegen schließen, wenn die Erhöhung auf 15 Thaler, und mit einem Minus von 3448 Thalern, wenn sie auf 20 Thaler festgesetzt würde. Ein so geringes Minus wie das im ersten Fall würde sich aus den eben angeführten Gründen schon nach kurzer Zeit ausgleichen. Die Festsetzung des variablen Teils der Pension auf 65 Thaler (also bei mehr als 18 Witwen die Verteilung der Summe von 1170 Thalern) wäre daher an sich überhaupt nicht bedenklich. Vielleicht würde man aber nicht gern von dem bisher stets beachteten Brauch abweichen, nach dem die halbjährliche Pensionssumme immer eine ganze Anzahl Pistolen betragen hat (wobei diese kleine Bequemlichkeit allerdings von selbst wegfallen wird, sobald die Zahl der Witwen über 18 gestiegen ist). Den variablen Teil aber schon jetzt auf 70 Thaler zu setzen, hielte ich schon aus \SourceLeaf{174}Prinzip nicht für ratsam, auch wenn ich unter den jetzigen günstigen Umständen gern die Hoffnung teile, dass sich das Minus von 3448 Thalern schon in den kommenden Jahren bedeutend vermindern würde.

\medskip
Ich glaube, der Universitätskirchendeputation im Vorstehenden ausreichend Material zusammengetragen zu haben, auf dessen Grundlage mit bewusster Sicherheit ein Beschluss gefasst werden kann. Ich bin aber gern zu weiteren Erläuterungen bereit, wenn diese für nötig gehalten werden sollten.

Göttingen, den 19. Oktober 1851.\hfill C. F. \textsc{Gauss}.
\begin{center}\rule{.28\linewidth}{.3pt}\end{center}
